\section{Conclusões}
\label{sec:conclusao}

O artigo examinou a associação dos preços reais de soja e milho com a conversão bruta da vegetação nativa para agricultura em Mato Grosso, considerando a capacidade produtiva potencial, as perdas anteriores e condicionantes territoriais. A distinção entre produzir mais e abrir novas áreas orientou tanto a descrição quanto a comparação com pastagens. O painel combinou informações de transição, produção, potencial agronômico, preços e clima, preservando a diferença entre capacidade, desempenho realizado e resultado econômico.

A exigência de histórico para construir perdas alterou a composição da amostra e reduziu os gradientes preço×capacidade antes da inclusão dos termos de H1b. No núcleo, os sinais de H1a permaneceram positivos, mas imprecisos; os modificadores associados às perdas também apresentaram intervalos que abrangem zero nas três inferências. Esses resultados não confirmam as sub-hipóteses de forma geral nem demonstram ausência das associações propostas. As comparações por período, destino e exposição acrescentaram heterogeneidade, inclusive resultado contrário à direção de H1a, cuja interpretação depende do suporte e do tratamento da dependência.

As perdas físicas não identificam, isoladamente, restrições financeiras ou intenção de compensar receita por expansão. Da mesma forma, frações territoriais e datas institucionais não medem cumprimento nem permitem atribuir os contrastes ao Código Florestal ou à Moratória. A conexão entre ambiente, condições locais e sequência temporal permanece pertinente, com o alcance empírico delimitado pelas variáveis observadas. Uma defasagem organiza a comparação, sem demonstrar planejamento intertemporal.

Para a gestão ambiental, o estudo reforça a necessidade de acompanhar destinos da abertura, condições produtivas e cobertura da informação de forma articulada. O crescimento da produção não basta para inferir pressão de conversão, e uma associação municipal não define a resposta de cada produtor. Avançar na ligação entre risco agrícola e conservação requer observar a transmissão das perdas aos recursos e às decisões sobre a terra, além das obrigações e dos compromissos por imóvel ou cadeia. Essas informações permitiriam avaliar instrumentos específicos, preservando as diferenças entre os biomas e as alternativas efetivamente acessíveis aos produtores.
